\documentstyle[% 


righttag% 


,11pt% 


,amssymb% 


,russian%               remove in Eng. 


,izvru%                 Set izven% in Eng. 


]{amsart} 


 


%  


% verify \bold 0 


%% \renewcommand{\baselinestretch}{0.98}


\newcommand{\la}{\langle} 


\newcommand{\ra}{\rangle} 


\newcommand{\myfont}{\series{m}\shape{n}\selectfont} 


\newcommand{\SP}{\operatorname{Sp}} 


\newcommand{\up}[2]{\overset{#1}{#2}{}} 


\newcommand{\tts}[1]{} 


 


\theoremstyle{plain} 


\theorembodyfont{\it} 


\newtheorem{thms}{\hskip\parindent Теорема}[section] 


\newtheorem{lems}{\hskip\parindent Лемма}[section] 


 


\theoremstyle{definition} 


\theorembodyfont{\rm} 


\newtheorem{defns}{\hskip\parindent Определение}[section] 


 


\theoremstyle{definition} 


\theorembodyfont{\rm} 


\newtheorem{cornonum}{\hskip\parindent Следствие} 


\newtheorem{notenonum}{\hskip\parindent Замечание} 


\renewcommand{\thecornonum}{} 


\renewcommand{\thenotenonum}{} 


\numberwithin{equation}{section} 


 


%\hoffset=-15mm%                  For Eng. 


\setcounter{page}{34} 


\god{2003} 


\nomer{3~(490)} %                        Remove in Eng. 


%\nomer{Vol.\,47, No.\,3}%               Set in Eng. 


%\str{\pageref{begin}--\pageref{end}}%  Set in Eng. 


\udk{517.977} 


 


\author{\it Ю.Н.\,Панкратова} 


% NB:   ^^ remove in Eng. 


\title{Равновесие по Нэшу в смешанном расширении игры\\ 


с упорядоченными исходами} 


 


\date{} 


%\pagestyle{myheadings}%         Set in Eng. 


\pagestyle{plain} %              Remove in Eng. 


\begin{document} 


%\thispagestyle{plain}%          Set in Eng. 


\maketitle 


\label{begin} 


 


 


Статья посвящена проблеме существования и описания ситуаций 


равновесия по Нэшу в смешанном расширении конечной игры с 


упорядоченными исходами. Как известно, для игр с числовыми выигрышами 


принцип равновесия в смысле Нэша реализуется в смешанных стратегиях. Для 


игр с упорядоченными исходами конструкция смешанного расширения может 


быть введена следующим образом: множества стратегий и исходов 


расширяются до множеств вероятностных мер, а функции реализации и 


отношения порядка игроков продолжаются на множества вероятностных мер. 


 


При построении смешанного расширения игры с упорядоченными 


исходами основной проблемой является нахождение способа продолжения 


порядка на множество вероятностных мер. В данной работе используется так 


называемое коническое продолжение порядка на множество вероятностных 


мер, которое задается при помощи конусов изотонных отображений 


упорядоченных множеств в $R$. В [1] показано, что в смешанном расширении 


игры с упорядоченными исходами, построенном при помощи конечнопорожденных 


конусов, существуют ситуации равновесия. Однако наиболее важного типа 


равновесия, равновесия по Нэшу, в таком расширении может не быть. 


 


Цель данной работы --- нахождение необходимых, а также достаточных 


условий существования ситуаций равновесия по Нэшу в смешанном 


расширении игры с упорядоченными исходами. Основной метод решения 


данной задачи состоит в характеризации системы спектров равновесной 


ситуации, которая базируется на условии дополняющей нежесткости (заметим, 


что для игр с упорядоченными исходами условие дополняющей нежесткости 


является гораздо более сильным условием, чем для игр с числовыми 


функциями выигрыша). 


 


В первой части статьи находятся условия, при которых смешанное 


расширение игры с упорядоченными исходами, построенное с помощью 


конусов изотонных отображений, будет игрой с упорядоченными исходами. 


 


Во второй части дается характеризация спектров ситуации равновесия по 


Нэшу в смешанном расширении игры с упорядоченными исходами, причем 


основным здесь является условие сбалансированности. 


 


Для игры двух лиц с упорядоченными исходами условие вполне 


равновесности ситуации в смешанном расширении сводится к условию 


сбалансированности матрицы ее исходов [2], которое в свою очередь 


эквивалентно коллинеарной сбалансированности некоторых покрытий 


множеств стратегий игроков [3]. В третьей части работы этот результат 


обобщается на игры $n$ лиц, причем к условию коллинеарности добавляется 


еще одно дополнительное условие. Отметим, что понятие 


сбалансированности 


покрытий введено в [4], однако до сих пор оно использовалось в теории игр 


только применительно к семействам коалиций игроков в рамках теории 


кооперативных игр. 


\pagebreak


 


\section{Расширение игры с упорядоченными исходами\protect\\ с помощью конусов 


изотонных отображений} 


 


Стратегическая игра $n$ лиц с упорядоченными исходами формально 


может быть задана в виде набора объектов 


\begin{equation} 


G=\la I,(X_i)_{i\in I},A,(\omega_i)_{i\in I},F\ra, 


\end{equation} 


где $I=\{1,2,\dotsc,n\}$ --- множество игроков; $X_i$ --- множество 


стратегий (или чистых стратегий) игрока $i\in I$; $A$ --- множество 


исходов; $\omega_i$ --- отношение (частичного) 


порядка на множестве $A$, выражающее предпочтения игрока $i$; 


$F:\prod\limits_{i\in I}X_i\to A$ --- 


функция реализации, которая каждой ситуации игры $x\in 


X=\prod\limits_{i\in I}X_i$ ставит в 


соответствие определяемый ею исход $F(x)\in A$. 


 


Смешанное расширение игры $G$ с упорядоченными исходами 


представляет собой игру, в которой базисные множества заменяются 


множествами вероятностных мер, а функция реализации и отношения порядка, 


выражающие предпочтения игроков, продолжаются на множества 


вероятностных мер. 


 


Далее мы ограничиваемся конечными играми. В этом случае смешанное 


расширение игры $G$ есть игра $\wt G$, в которой множеством стратегий 


игрока $i$ является множество $\wt X_i$ вероятностных векторов над 


$X_i$, а множеством исходов 


--- множество $\wt A$ вероятностных векторов над $A$. Функция реализации в 


смешанном расширении игры $G$ есть отображение $\wt F$, которое каждой 


ситуации в смешанных стратегиях $(\mu_1,\mu_2,\dotsc,\mu_n)\in 


\prod\limits_{i\in I}\wt X_i$ ставит в соответствие 


вероятностный вектор 


\begin{equation} 


\wt F_{(\mu_1,\mu_2,\dotsc,\mu_n)}(a)=\sum\begin{Sb} 


F(x_1,x_2,\dotsc,x_n)=a \\ (x_1,x_2,\dotsc,x_n)\in X\end{Sb} 


\mu_1(x_1)\mu_2(x_2)\dotsc\mu_n(x_n). 


\end{equation} 


 


Продолжение порядка $\omega_i$ на множество вероятностных мер $\wt A$ 


строится следующим образом [1]. Для каждого $i\in I$ зафиксируем 


некоторый конус $C_i$ изотонных отображений упорядоченного множества 


$\la A,\omega_i\ra$ в $R$. Положим для произвольных $\mu,\nu\in A$ 


\begin{equation} 


\mu\up{\wt \omega_i(C_i)}{\le}\nu\ \Leftrightarrow 


\ (\varphi,\mu)\le(\varphi,\nu) \ \; \forall\varphi\in C_i, 


\end{equation} 


где $(\ ,\ )$ --- стандартное скалярное произведение в $R^A$. Построенная 


игра обозначается\linebreak $\wt G[C_1,C_2,\dotsc,C_n]$. 


 


Первая задача состоит в нахождении условий, накладываемых на конусы 


изотонных отображений $(C_i)_{i\in I}$, при которых расширенная игра 


$\wt G[C_1,C_2,\dotsc,C_n]$ 


оказывается игрой с упорядоченными исходами, а $G$ --- ее подыгрой. 


 


\begin{defns} 


Игра $\ov G=\la I,(\ov X_i)_{i\in I},\ov A,(\ov \omega_i)_{i\in I},\ov 


F\ra$ есть расширение игры $G$, если при каждом $i\in I$ 


$$ 


X_i\subseteq\ov X_i,\ \ A\subseteq\ov A, 


\ \ \omega_i\subseteq\ov\omega_i,\ \ F\subseteq\ov F. 


$$ 


\end{defns} 


 


\begin{defns} 


Игра $G$ есть подыгра игры $\ov G$, если $\ov G$ есть расширение игры 


$G$ и 


\begin{alignat*}{2} 


&\text{а)} &\ \; \ov\omega_i|_A&=\omega_i,\ \ i\in I;\\ 


&\text{б)} &\ \; \ov F|_X&=F. 


\end{alignat*} 


\end{defns} 


 


\begin{notenonum} 


При рассмотрении смешанного расширения мы отождествляем, как 


обычно, вырожденный вероятностный вектор $\delta_z$, сосредоточенный в 


точке $z$, с этой точкой, т.\,е. $\delta_z\equiv z$. 


\end{notenonum} 


 


\begin{thms} 


Предположим, что каждый конус $C_i$ содержит ненулевой вектор с 


равными компонентами. Для того чтобы игра $\wt G[C_1,C_2,\dotsc,C_n]$ была 


игрой с 


упорядоченными исходами, а $G$ --- ее подыгрой, необходимо и достаточно, 


чтобы для каждого $i\in I$ 


 


$1)$ конус $C_i$ аппроксимировал отношение порядка $\omega_i;$ 


 


$2)$ конус $C_i$ был телесным в линейном пространстве $R^A$. 


\end{thms} 


 


Доказательство теоремы сводится к следующим леммам. 


 


\begin{lems} 


При любом наборе конусов $C_1,C_2,\dotsc,C_n$ игра $\wt 


G[C_1,C_2,\dotsc,C_n]$ 


является расширением игры $G$, т.\,е. 


$$ 


X_i\subseteq\wt X_i,\ \ A\subseteq\wt A, 


\ \ \omega_i\subseteq\wt\omega_i,\ \ F\subseteq\wt F. 


$$ 


\end{lems} 


 


\begin{pf} 


При отождествлении вырожденного вероятностного 


вектора, сосредоточенного в одной точке, с этой точкой получим включения 


$X_i\subseteq\wt X_i$, $A\subseteq\wt A$. С помощью (1.2) легко 


проверяется равенство 


$$ 


\wt F_{(\delta_{x_1},\delta_{x_2},\dotsc,\delta_{x_n})} 


=\delta_{F(x_1,x_2,\dotsc,x_n)}. 


$$ 


С учетом указанного выше отождествления имеем $\wt F|_X=F$, тем более 


$F\subseteq\wt F$. 


Докажем включение $\omega_i\subseteq\wt\omega_i(C_i)$. Рассмотрим 


произвольные элементы $a,a'\in A$ 


такие, что $(a,a')\in\omega_i$. Считая $A=\{a_1,a_2,\dotsc,a_m\}$ и 


рассматривая $\varphi$ и $\delta_{a_k}$ как 


$m$-компонентные векторы, имеем 


\begin{equation} 


(\varphi,\delta_{a_k})=\varphi(a_k). 


\end{equation} 


Так как всякое отображение $\varphi\in C_i$ изотонно относительно 


порядка $\omega_i$, при произвольном $\varphi\in C_i$ получим 


$$ 


a\up{\omega_i}{\le}a'\ \Rightarrow\ \varphi(a)\le\varphi(a') 


\ \up{(1.4)}{\Leftrightarrow}\ (\varphi,\delta_a)\le 


(\varphi,\delta_{a'})\ \up{\text{в силу (1.3)}}{\Leftrightarrow} 


\ \delta_a\up{\wt\omega_i(C_i)}{\le}\delta_{a'}. 


$$ 


С учетом указанного отождествления получим 


$\omega_i\subseteq\wt\omega_i(C_i)$. Таким образом, 


игра $\wt G$ является расширением игры $G$. 


\end{pf} 


 


\begin{lems} 


Игра $G$ является подыгрой игры $\wt G[C_1,C_2,\dotsc,C_n]$ тогда и только 


тогда, когда конусы $(C_i)_{i\in I}$ аппроксимирующие. 


\end{lems} 


 


\begin{pf} 


Игра $G$ будет подыгрой игры $\wt G[C_1,C_2,\dotsc,C_n]$ тогда и только 


тогда, когда $\wt\omega_i(C_i)|_A=\omega_i$. Поскольку согласно лемме 


1.1 \ $\omega_i\subseteq\wt\omega_i(C_i)$, то 


$\omega_i\subseteq\wt\omega_i(C_i)\cap A^2=\wt\omega_i(C_i)|_A$. Обратное 


включение $\wt\omega_i|_A\subseteq\omega_i$ сводится к импликации 


$$ 


(\delta_a,\delta_{a'})\in\wt\omega_i(C_i)\ \Rightarrow 


\ (a,a')\in\omega_i. 


$$ 


 


По построению продолжения $\wt\omega_i(C_i)$ и с учетом (1.4) условие 


$(\delta_a,\delta_{a'})\in\wt\omega_i(C_i)$ эквивалентно условию 


$\varphi(a)\le\varphi(a')$ при произвольном $\varphi\in C_i$. 


Таким образом, игра $G$ будет подыгрой игры $\wt G[C_1,C_2,\dotsc,C_n]$ 


тогда и только 


тогда, когда условие $\varphi(a)\le\varphi(a')$\; $\forall\varphi\in 


C_i$ влечет $a\up{\omega_i}{\le}a'$, а это есть определение 


аппроксимируемости отношения порядка $\omega_i$ конусом $C_i$. 


\end{pf} 


 


\begin{lems} 


Если конусы $(C_i)_{i\in I}$ телесные, то игра $\wt 


G[C_1,C_2,\dotsc,C_n]$ является игрой с упорядоченными исходами. 


\end{lems} 


 


\begin{pf} 


В общем случае игра $\wt G[C_1,C_2,\dotsc,C_n]$ является игрой с 


квазиупорядоченными исходами. Покажем, что в нашем случае отношения 


квазипорядка $\wt\omega_i(C_i)$ будут антисимметричными, т.\,е. они будут 


отношениями 


порядка. Надо проверить выполнимость импликации 


\begin{equation} 


\mu\up{\wt\omega_i(C_i)}{\le}\nu, 


\ \ \nu\up{\wt\omega_i(C_i)}{\le}\mu\ \Rightarrow\ \mu=\nu. 


\end{equation} 


 


В силу (1.3) условие импликации (1.5) сводится к тому, что 


\begin{equation} 


(\varphi,\mu-\nu)=0\ \; \forall\varphi\in C_i.


\end{equation} 


Так как конус $C_i$ телесный, то существует $m=|A|$ линейно независимых 


векторов $p^1,p^2,...,p^m{\in} C_i$. С учетом (1.6) для любого 


%%%%% return dotsc ^^


$l=1,2,\dotsc,m$ выполняется 


$$ 


(p^l,\mu-\nu)=0. 


$$ 


Таким образом, координаты вектора $\mu-\nu$ удовлетворяют однородной 


системе 


$m$ линейных уравнений с $m$ неизвестными. В силу линейной независимости 


векторов $p^1,p^2,\dotsc,p^m$ определитель этой системы отличен от нуля, 


следовательно, система имеет только тривиальное решение $\mu-\nu=\bold 0$, 


т.\,е. $\mu=\nu$. 


\end{pf} 


 


\begin{notenonum} 


До сих пор мы не использовали условие существования для каждого 


$C_i$ ненулевого вектора с одинаковыми компонентами. При выполнении этого 


условия верно утверждение, обратное лемме 1.3. 


\end{notenonum} 


 


{\bf Обратное утверждение.} {\it Если $\wt\omega_i(C_i)$ --- отношение 


порядка и $C_i$ содержит 


ненулевой вектор с одинаковыми компонентами, то конус $C_i$ телесный.} 


\medskip 


 


{\bf Доказательство.} Предположим, что $\dim C_i=r<m=|A|$, тогда 


максимальная линейно независимая подсистема в $C_i$ содержит $r$ 


векторов. Пусть $q^1,q^2,\dotsc,q^r$ --- максимальная линейно 


независимая система векторов в $C_i$. Составим систему 


$r$ уравнений относительно $y=(y_1,y_2,\dotsc,y_m)$ 


\begin{equation} 


(q^j,y)=0,\quad j=1,\dotsc,r. 


\end{equation} 


 


Так как ранг матрицы системы (1.7) равен $r<m$, система имеет 


нетривиальное решение $y^*=(y^*_1,y^*_2,\dotsc,y^*_m)\ne\bold 0$. 


По предположению конус $C_i$ содержит ненулевой вектор с одинаковыми 


компонентами $\varphi_0=(k,k,\dotsc,k)\in C_i$, где $k\ne0$, причем 


$\varphi_0$ может быть линейно 


выражен через векторы $q^1,q^2,\dotsc,q^r:\varphi_0= 


\sum\limits^m_{j=1}\lambda_j q^j$. В силу (1.7) 


$$ 


(\varphi_0,y^*)=\sum^m_{j=1}\lambda_j(q^j,y^*)=0; 


$$ 


с другой стороны, $(\varphi_0,y^*)=k y^*_1+k y^*_2+\dotsb+k 


y^*_m=k(y^*_1+y^*_2+\dotsb+y^*_m)$,


следовательно, 


\begin{equation} 


y^*_1+y^*_2+\dotsb+y^*_m=0. 


\end{equation} 


 


Зафиксируем вероятностный вектор $u=(u_1,u_2,\dotsc,u_m)\in\wt A$, где 


$0<u_i<1$, $\sum\limits^m_{j=1}u_j=1$. Рассмотрим вектор $v=u+\alpha 


y^*$, причем возьмем $\alpha>0$ 


настолько малым, чтобы $0<v_j<1$, $j=1,2,\dotsc,m$. В силу (1.8) 


$$ 


\sum^m_{j=1}v_j=\sum^m_{j=1}(u_j+\alpha y^*_j)= 


\sum^m_{j=1}u_j+\alpha\sum^m_{j=1}y^*_j=\sum^m_{j=1}u_j=1, 


$$ 


таким образом, $v\in\wt A$. 


 


Так как любой вектор $\varphi\in C_i$ может быть представлен в виде 


$\varphi=\sum\limits^m_{j=1}\beta_j q^j$, то 


$$ 


(\varphi,v-u)=(\varphi,\alpha y^*)=\alpha(\varphi,y^*)=\alpha\bigg( 


\sum^m_{j=1}\beta_j q^j,y^*\bigg)=\alpha\sum^m_{j=1} 


\beta_j(q^j,y^*)=0, 


$$ 


откуда $(\varphi,v)=(\varphi,u)$\; $\forall\varphi\in C_i$, т.\,е. 


$u\up{\wt\omega_i(C_i)}{\le}v$ и $v\up{\wt\omega_i(C_i)}{\le}u$. Так как 


$y^*\ne\bold 0$, то $u\ne v$. 


Следовательно, условие антисимметричности для квазипорядка 


$\wt\omega_i(C_i)$ не имеет места.


Теорема 1.1 доказана.


 


 


\section{Ситуации равновесия по Нэшу в игре $\wt G[C_1,C_2,\dotsc,C_n]$} 


 


Рассмотрим игру $\wt G[C_1,C_2,\dotsc,C_n]$, являющуюся смешанным 


расширением игры с упорядоченными исходами $G$ с помощью конусов 


$[C_1,C_2,\dotsc,C_n]$. Установим некоторые общие свойства игры $\wt 


G[C_1,C_2,\dotsc,C_n]$. 


 


Для $\mu=(\mu_1,\mu_2,\dotsc,\mu_n)\in\prod\limits_{i\in I}\wt X_i$ 


через $\mu\|\mu^*_i$ обозначим ситуацию, 


полученную заменой стратегии $\mu_i$ на стратегию $\mu^*_i$. 


 


\begin{defns} 


Ситуация $\mu^*=(\mu^*_1,\mu^*_2,\dotsc,\mu^*_n)\in \prod\limits_{i\in 


I}\wt X_i$ называется ситуацией 


равновесия по Нэшу в игре $\wt G[C_1,C_2,\dotsc,C_n]$, если для всех 


$\mu_i\in\wt X_i$, $i\in I$,  выполняется 


\begin{equation} 


\wt F_{(\mu^*\|\mu_i)}\up{\wt\omega_i(C_i)}{\le}\wt F_{\mu^*}. 


\end{equation} 


\end{defns} 


 


\begin{thms} 


Пусть $C_i$ при каждом $i\in I$ --- аппроксимирующий конус 


изотонных отображений упорядоченного множества $\la A,\omega_i\ra$ в $R$. 


Для того чтобы ситуация $\mu^*=(\mu^*_1,\mu^*_2,\dotsc,\mu^*_n)\in 


\prod\limits_{i\in I}\wt X_i$ была ситуацией равновесия по 


Нэшу в игре $\wt G[C_1,C_2,\dotsc,C_n]$, необходимо и достаточно, чтобы 


неравенства $(2.1)$ 


выполнялись при чистых отклонениях игроков, т.\,е. при всех 


$x_i\in X_i$, $i\in I$,  имело место 


\begin{equation} 


\wt F_{(\mu^*\|\delta_{x_i})}\up{\wt\omega_i(C_i)}{\le}\wt F_{\mu^*}. 


\end{equation} 


\end{thms} 


 


Доказательство теоремы основано на следующих двух несложных 


леммах, доказательство которых опускаем. 


 


\begin{lems} 


Функция реализации $\wt F$ линейна, т.\,е. для


$\sum\limits^k_{l=1}\alpha_l=1$, 


$\mu^1_i,\dotsc,\mu^k_i\in\wt X_i$ 


при каждом $i\in I$ и любых 


$\alpha_1,\alpha_2,\dotsc,\alpha_k\ge0$ выполняется 


$$ 


\wt 


F_{\left(\mu\left\|\sum\limits^k_{l=1}\right.\alpha_l\mu^l_i\right)} 


=\sum^k_{l=1}\alpha_l\wt F_{(\mu\|\mu^l_i)}. 


$$ 


\end{lems} 


 


\begin{cornonum}[{\myfont 


разложение функции реализации по чистым стратегиям}] 


Для функции реализации $\wt F$ справедливо разложение 


\begin{equation} 


\wt F_{(\mu_1,\mu_2,\dotsc,\mu_n)}=\sum_{x_i\in X_i} 


\mu_i(x_i)\wt F_{(\mu\|\delta_{x_i})}. 


\end{equation} 


\end{cornonum} 


 


Доказательство непосредственно вытекает из леммы 2.1 с учетом равенства 


$$ 


\mu_i=\sum_{x_i\in X_i}\mu_i(x_i)\delta_{x_i}. 


$$ 


 


{\bf Лемма 2.2}


(свойство стабильности отношения порядка $\wt\omega(C)$ относительно 


операции смешивания). {\it 


Пусть $C$ --- некоторый конус в линейном пространстве 


$R^A$, состоящий из изотонных отображений упорядоченного множества $\la 


A,\omega\ra$ в 


$R$, $\wt\omega(C)$ --- продолжение порядка $\omega$ на $A$, определяемое 


$(1.3)$, и $\alpha_l\ge0$, $\sum\limits^k_{l=1}\alpha_l=1$. 


Тогда из условий 


\begin{equation} 


\mu_i\up{\wt\omega(C)}{\le}\nu_i,\quad i=1,\dotsc,k, 


\end{equation} 


следует 


\begin{equation} 


\alpha_1\mu_1+\alpha_2\mu_2+\dotsb+\alpha_k\mu_k 


\up{\wt\omega(C)}{\le}\alpha_1\nu_1+\alpha_2\nu_2+\dotsb+ 


\alpha_k\nu_k. 


\end{equation} 


}


\addtocounter{lems}{1}


 


{\bf Доказательство теоремы 2.1. Необходимость} очевидна, т.\,к. чистая 


стратегия является частным случаем смешанной стратегии. 


\smallskip 


 


{\bf Достаточность.} Покажем, что выполнение неравенств (2.2) влечет 


выполнение (2.1). Рассмотрим смешанную стратегию $\mu_i\in\wt X_i$. 


Умножая обе части (2.2) на $\mu_i(x_i)$ 


и суммируя по $x_i\in X_i$, по лемме 2.2 получим 


$$ 


\sum_{x_i\in X_i}\mu_i(x_i)\wt F_{(\mu^*\|\delta_{x_i})} 


\up{\wt\omega_i(C_i)}{\le}\sum_{x_i\in X_i}\mu_i(x_i) 


\wt F_{\mu^*}. 


$$ 


Используя разложение функции реализации по чистым стратегиям (2.3), имеем 


$$ 


\sum_{x_i\in X_i}\mu_i(x_i)\wt F_{(\mu^*\|\delta_{x_i})} 


=\wt F_{(\mu^*\|\mu_i)}. 


$$ 


Так как $\sum\limits_{x_i\in X_i}\mu_i(x_i)\wt F_{\mu^*}=\wt 


F_{\mu^*}\sum\limits_{x_i\in X_i}\mu_i(x_i)=\wt F_{\mu^*}$, то в итоге 


получим $\wt F_{(\mu^*\|\mu_i)}\up{\wt\omega_i(C_i)}{\le}\wt 


F_{\mu^*}$.$\quad\square$ 


 


\begin{cornonum} 


В предположениях теоремы 2.1 произвольная ситуация в чистых 


стратегиях будет ситуацией равновесия по Нэшу в игре $\wt 


G[C_1,C_2,\dotsc,C_n]$ тогда и 


только тогда, когда она является ситуацией равновесия по Нэшу в игре $G$. 


\end{cornonum} 


 


Через $\SP\mu_i=\{x_i:\mu_i(x_i)>0\}$ обозначим спектр вероятностной меры 


$\mu_i$. 


 


Теорема 2.1 может быть уточнена в отношении свойств системы спектров 


ситуации равновесия по Нэшу следующим образом. 


 


\begin{thms}[{\myfont 


правило дополняющей нежесткости}] 


Пусть $C_i$ при каждом $i\in I$ --- 


аппроксимирующий телесный конус изотонных отображений 


упорядоченного множества $\la A,\omega_i\ra$ в $R$. Для того чтобы 


ситуация $\mu^*\in\prod\limits_{i\in I}\wt X_i$ 


была ситуацией равновесия по Нэшу в игре $\wt G[C_1,C_2,\dotsc,C_n]$, 


необходимо и достаточно, чтобы для любого $i\in I$ выполнялись условия 


\begin{equation} 


\begin{split} 


1)\ \; &\wt F_{(\mu^*\|\delta_{x'_i})}=\wt F_{\mu^*} 


\qquad \forall x'_i\in\SP\mu^*_i; \\ 


2)\ \; &\wt F_{(\mu^*\|\delta_{x''_i})}\up{\wt\omega_i(C_i)}{\le} 


\wt F_{\mu^*}\ \ \forall x''_i\notin\SP\mu^*_i. 


\end{split} 


\end{equation} 


\end{thms} 


 


\begin{pf} 


{\bf Достаточность.} Условия (2.6) есть частный случай условий 


(2.2). По теореме 2.1 \ $\mu^*$ --- ситуация равновесия по Нэшу в игре 


$\wt G[C_1,C_2,\dotsc,C_n]$. 


Доказательство необходимости основано на следующей лемме. 


 


\begin{lems}[{\myfont 


уточненное свойство стабильности}] 


Пусть $C$ --- некоторый конус 


в линейном пространстве $R^A$, состоящий из изотонных отображений 


упорядоченного множества $\la A,\omega\ra$ в $R$, $\wt\omega(C)$ --- 


продолжение порядка $\omega$ на $\wt A$, 


определяемое $(1.3)$, $\alpha_l>0$, $\sum\limits^k_{l=1}\alpha_l=1$. 


Тогда из условий


$$


\begin{cases}


\mu_i\up{\wt\omega(C)}{<}\nu_i, & i=1,\dotsc,k-1,\\


\mu_k\up{\wt\omega(C)}{\le}\nu_k & \


\end{cases}


$$ 


следует


\begin{equation} 


\alpha_1\mu_1+\alpha_2\mu_2+\dotsb+


\alpha_k\mu_k \up{\wt\omega(C)}{<}


\alpha_1\nu_1+\alpha_2\nu_2+\dotsb+\alpha_k\nu_k.


\end{equation}


\end{lems} 


 


Докажем {\bf необходимость} в теореме 2.2. Пусть $\mu^*$ --- ситуация 


равновесия по Нэшу. Тогда по теореме 2.1 выполняются неравенства (2.2). 


Для доказательства теоремы осталось показать, что для любого $i\in I$ 


для чистых стратегий, принадлежащих спектру $\mu^*_i$, неравенства 


(2.2) обращаются в равенства. 


Предположим противное, т.\,е. что для некоторого $x'_i\in\SP\mu^*_i$ имеет 


место $\wt F_{(\mu^*\|\delta_{x'_i})}\ne\wt F_\mu^*$. Так как по лемме 


1.3 \ $\wt\omega(C_i)$ --- отношение порядка, 


то выполняется строгое неравенство 


$$ 


\wt F_{(\mu^*\|\delta_{x'_i})}\up{\wt\omega_i(C_i)}{<}\wt F_{\mu^*}. 


$$ 


Таким образом, система (2.2) содержит строгое неравенство. Умножая каждое 


неравенство системы (2.2) на $\mu^*_i(x_i)>0$, где $x_i\in\SP\mu^*_i$, и 


суммируя по $x_i\in\SP\mu^*_i$, по лемме 2.3 получаем 


\begin{equation} 


\sum_{x_i\in\SP\mu_i}\mu^*_i(x_i) 


\wt F_{(\mu^*\|\delta_{x_i})}\up{\wt\omega_i(C_i)}{<} 


\sum_{x_i\in\SP\mu_i}\mu^*_i(x_i)\wt F_{\mu^*}. 


\end{equation} 


 


Рассмотрим левую часть (2.8) 


\begin{multline*} 


\sum_{x_i\in\SP\mu_i^*}\mu^*_i(x_i) 


\wt F_{(\mu^*\|\delta_{x_i})}=\sum_{x_i\in\SP\mu_i^*} 


\mu^*_i(x_i)\wt F_{(\mu^*\|\delta_{x_i})}+ 


\sum_{x_i\notin\SP\mu_i^*}\mu^*_i(x_i)


\wt F_{(\mu^*\|\delta_{x_i})}= \\ 


%


=\sum_{x_i\in X_i}


\mu^*_i(x_i)\wt F_{(\mu^*\|\delta_{x_i})}\up{(2.4)}{=} 


\wt F_{(\mu^*)}. 


\end{multline*} 


Поскольку правая часть (2.8) есть $\wt F_{\mu^*}$, то


%% $$ 


%% \sum_{x_i\in\SP\mu_i^*}\mu^*_i(x_i)\wt F_{\mu^*}=\wt F_{\mu^*}. 


%% $$ 


в итоге из (2.8) получим $\wt F_{(\mu^*)}<\wt F_{\mu^*}$ --- противоречие. 


 


Таким образом, при чистых отклонениях игрока $i$ в пределах спектра 


стратегии $\mu^*_i$ выполняются неравенства (2.7). 


\end{pf} 


 


 


\section{Ситуации равновесия и сбалансированные покрытия\protect\\ в играх с 


упорядоченными исходами} 


 


Для игры двух лиц с упорядоченными исходами условие вполне 


равновесности ситуации в смешанном расширении сводится к условию 


сбалансированности матрицы ее исходов, которое в свою очередь 


эквивалентно коллинеарной сбалансированности некоторых покрытий 


множеств стратегий игроков [3]. Здесь мы обобщаем этот результат на игры 


$n$ лиц, причем к условию коллинеарности добавляется еще одно 


дополнительное условие, которое в случае игры двух лиц выполняется 


автоматически. 


 


\begin{defns} 


Матрица $\|F(x)\|_{x\in X}$ исходов игры $G$ вида (1.1) называется 


{\it сбалансированной}, если существует такая ситуация 


$\mu=(\mu_1,\mu_2,\dotsc,\mu_n)\in \prod\limits_{i\in I}\wt X_i$ с 


положительными компонентами, что при произвольном $a\in A$ и любых 


$x_{i_1}\in X_{i_1}$, $x_{i_2}\in X_{i_2}$, $i_1,i_2\in I$, выполняется 


$$ 


\wt F_{(\mu\|\delta_{x_{i_1}})}(a)= 


\wt F_{(\mu\|\delta_{x_{i_2}})}(a). 


$$ 


При этом набор векторов $\mu_1,\mu_2,\dotsc,\mu_n$ называется 


{\it балансовым набором} данной матрицы. 


\end{defns} 


 


\begin{defns} 


Пусть $E$ --- произвольное множество. Покрытие $(E_1,E_2,\dotsc,E_k)$ 


множества $E$ называется {\it сбалансированным}, если существует такой 


вектор $\lambda=(\lambda_1,\lambda_2,\dotsc,\lambda_k)$ с положительными 


компонентами, что для любого $e\in E$ имеет место 


$\sum\limits_{E_s\ni e}\lambda_s=1$. 


Вектор $\lambda$ называется {\it репрезентативным вектором} данного 


покрытия. 


\end{defns} 


 


Для каждого игрока $i\in I$ положим $X^i=\prod\limits_{j\in I\setminus 


i}X_j$ --- {\it множество наборов стратегий} всех остальных игроков. 


Так как пара $(x_i,y)\in X_i\times X^i$ определяет единственную 


ситуацию в игре $G$, функция реализации $F$ может быть 


представлена в виде $F(x_i,y)$. Для произвольных $a\in A$, $i\in I$, 


$y\in X^i$ рассмотрим семейство {\it характеристических} множеств 


$(S^a_y)_{y\in X^i}$, где подмножество $S^a_y\subseteq X_i$ 


определено следующим образом: 


$$ 


x_i\in S^a_y\ \Leftrightarrow\ F(x_i,y)=a. 


$$ 


 


\begin{thms} 


Для того чтобы матрица $\|F(x)\|_{x\in X}$ исходов игры $G$ была 


сбалансированной необходимо и достаточно, чтобы при любых $a\in A$, 


$i\in I$ семейство множеств $(S^a_y)_{y\in X^i}$ было сбалансированным 


покрытием множества $X_i$, имеющим репрезентативный вектор 


$(\lambda^a_y)_{y\in X^i}$, удовлетворяющий 


следующим условиям\/${:}$ 


 


$1)$ при любых $a,a'\in A$ отношение $\frac{\lambda^a_y}{\lambda^{a'}_y}$ 


есть константа по $y\in X^i;$ 


 


$2)$ при любых $a\in A$, $i\in I$, $i_0\in I\setminus i$, $y',y''\in 


X^i$ отношение $\frac{\lambda^a_{y'\|x_{i_0}}} 


{\lambda^a_{y''\|x_{i_0}}}$ есть 


константа по $x_{i_0}\in X_{i_0}$. 


\end{thms} 


 


\begin{pf} 


{\bf Необходимость.} Пусть $(\mu_1,\mu_2,\dotsc,\mu_n)$ --- балансовый 


набор для заданной матрицы. Положим $\gamma(a)=\wt 


F_{(\mu\|\delta_{x_1})}(a)=\wt F_{(\mu\|\delta_{x_2})}(a)=\dotsb= 


\wt F_{(\mu\|\delta_{x_n})}(a)$. 


При фиксированных $a\in A$, $i\in I$ для любого $y=(x_j)_{j\in 


I\setminus i}\in X^i$ определим 


$$ 


\lambda^a_y=\frac{\prod\limits_{j\in I\setminus i}\mu_j(x_j)} 


{\gamma(a)}>0. 


$$ 


Покажем, что вектор $(\lambda^a_y)_{y\in X^i}$ является репрезентативным 


вектором сбалансированного покрытия $(S^a_y)_{y\in X^i}$. Действительно, 


при произвольных фиксированных $x_i\in X_i$, $i\in I$, имеем 


$$ 


\sum_{S^a_y\ni x_i}\lambda^a_y= 


\sum\begin{Sb}F(x_i,y)=a \\ y\in X^i\end{Sb}\lambda^a_y 


=\sum\begin{Sb}F(x_i,y)=a \\ y\in X^i\end{Sb} 


\frac{\prod\limits_{j\in I\setminus i}\mu_j(x_j)}{\gamma(a)}= 


\frac{1}{\gamma(a)} 


\sum\begin{Sb}F(x_i,y)=a \\ y\in X^i\end{Sb} 


\prod_{j\in I\setminus i}\mu_j(x_j)=\frac{1}{\gamma(a)} 


\wt F_{(\mu\|\delta_{x_i})}(a)=1. 


$$ 


 


Далее проверим выполнение условий 1) и 2). 


 


1) При любых $a,a'\in A$ отношение $\frac{\lambda^a_y}{\lambda^{a'}_y}= 


\frac{\gamma(a')}{\gamma(a)}$ --- константа по $y\in X^i$. 


 


2) При любых $a\in A$, $i\in I$, $i_0\in I\setminus i$, 


$y'=(x'_j)_{j\in I\setminus i}$, $y''=(x''_j)_{j\in I\setminus i}$ 


имеем 


$$ 


\frac{\lambda^a_{y'\|x_{i_0}}}{\lambda^a_{y''\|x_{i_0}}}= 


\frac{\prod\limits_{j\in I\setminus\{i,i_0\}}\mu_j(x'_j) 


\mu_{i_0}(x_{i_0})}{\gamma(a)}\ :\ 


\frac{\prod\limits_{j\in I\setminus\{i,i_0\}}\mu_j(x''_j) 


\mu_{i_0}(x_{i_0})}{\gamma(a)}= 


\frac{\prod\limits_{j\in I\setminus\{i,i_0\}}\mu_j(x'_j)} 


{\prod\limits_{j\in I\setminus\{i,i_0\}}\mu_j(x''_j)} 


$$ 


не зависит от выбора $x_{i_0}\in X_{i_0}$, т.\,е. является константой по 


$x_{i_0}\in X_{i_0}$. Необходимость доказана. 


\medskip 


 


{\bf Достаточность.} Зафиксируем $j\in I$. Для каждого $x_j\in X_j$ 


определим $\mu_j(x_j)$ следующим образом. Возьмем $i\ne j$, $a'\in A$. 


Выберем $y'\in X^i$, для которого $y'_j=x_j$, тогда $y'=y'\|x_j$. 


Положим 


$$ 


\sigma_j(x_j)=(\lambda^{a'}_{y'\|x_j})^{1/(n-1)},\quad 


\mu_j(x_j)=\frac{\sigma_j(x_j)} 


{\sum\limits_{x^*_j\in X_j}\sigma_j(x^*_j)}. 


$$ 


Корректность определения величин $\mu_j(x_j)$, т.\,е. их независимость от 


выбора 


$a'\in A$ и $y'\in X^i$ следует из условий 1) и 2) теоремы. Ясно, что 


вектор $(\mu_j(x_j))_{x_j\in X_j}$ является вероятностным вектором над 


$X_j$, $j=1,2,\dotsc,n$. 


 


Проверим, что построенный набор $(\mu_1,\mu_2,\dotsc,\mu_n)$ является 


балансовым набором для матрицы $\|F(x)\|_{x\in X}$. Действительно, при 


любых $x_i\in X_i$, $i\in I$, $a\in A$ имеем 


$$ 


\wt F_{(\mu\|\delta_{x_i})}(a)=\sum\begin{Sb} 


F(x_i,y)=a \\ y\in X^i\end{Sb}\prod_{j\in I\setminus i} 


\mu_j(x_j)=\sum\begin{Sb}F(x_i,y)=a \\ y\in X^i \end{Sb} 


\prod_{j\in I\setminus i}\frac{\sigma_j(x_j)} 


{\sum\limits_{x^*_j\in X_j}\sigma_j(x^*_j)}= 


\sum\begin{Sb}F(x_i,y)=a\\ y\in X^i \end{Sb} 


\prod_{j\in I\setminus i} 


\frac{(\lambda^{a'}_{y'\|x_j})^{1/(n-1)}}


{\sum\limits_{x^*_j\in X_j} 


(\lambda^{a'}_{y'\|x^*_j})^{1/(n-1)}}. 


$$ 


Так как $\mu_j(x_j)$ не зависят от выбора $a'\in A$, положим $a'=a$. Выбор 


вектора $y'$ по построению зависит от $x_j$; в нашем случае при всех 


$j\in I\setminus i$ можно взять единый вектор $y'=y=(x_j)_{j\in J}$. 


Обозначим $\rho^i(a)=\prod\limits_{j\in I\setminus i} 


\sum\limits_{x^*_j\in X_j}(\lambda^a_{y\|x^*_j})^{1/(n-1)}$. Имеем 


\begin{equation} 


\wt F_{(\mu\|\delta_{x_i})}(a)=\sum 


\begin{Sb}F(x_i,y)=a \\ y\in X^i \end{Sb}\frac{1}{\rho^i(a)} 


\prod_{j\in I\setminus i}(\lambda^a_{y\|x_j})^{1/(n-1)} 


=\frac{1}{\rho^i(a)}\sum\begin{Sb} 


S^a_y\ni x_i \\ y\in X^i\end{Sb} 


\lambda^a_{y\|x_j}=\frac{1}{\rho^i(a)}. 


\end{equation} 


Показали, что $\wt F_{(\mu\|\delta_{x_i})}(a)$ не зависит от $x_i\in 


X_i$ при любом $i\in I$. Используя 


разложение функции реализации по чистым стратегиям (лемма 2.1) и (3.1), 


получим 


$$ 


\wt F_{(\mu)}(a)=\sum_{x_i\in X_i}\mu_i(x_i) 


\wt F_{(\mu\|\delta_{x_i})}(a)=\frac{1}{\rho^i(a)} 


\wt F_{(\mu\|\delta_{x_i})}(a). 


$$ 


Последнее равенство означает, что матрица $\|F(x)\|_{x\in X}$ является 


сбалансированной, а набор $(\mu_1,\mu_2,\dotsc,\mu_n)$ является ее 


балансовым набором. 


\end{pf}


 


\begin{cornonum} 


Для того чтобы игра $\wt G[C_1,C_2,\dotsc,C_n]$ 


имела ситуацию вполне равновесия в смысле Нэша, необходимо и 


достаточно, чтобы матрица исходов игры $G$ удовлетворяла условиям 


теоремы 3.1. 


\end{cornonum} 


 


Таким образом, в условиях теоремы 2.1 существование ситуации вполне 


равновесия по Нэшу в коническом расширении игры $G$ определяется 


исключительно функцией реализации игры $G$. 


 


\begin{thebibliography}{9} 


 


\bibitem{1} 


Розен В.В. {\em Описание ситуаций равновесия в играх с упорядоченными 


исходами} // Математика, механика, кибернетика. --


Саратов: Изд-во СГУ, 1999. -- С.\,128--131. 


\bibitem{2} 


Сердюкова Ю.Н. {\em Ситуации равновесия по Нэшу с заданными 


спектрами} // Математика, механика, кибернетика. --


Саратов: Изд-во СГУ, 1999. -- С.\,137--139. 


\bibitem{3} 


Розен В.В., Панкратова Ю.Н. {\em Ситуации равновесия и сбалансированные 


покрытия в играх с упорядоченными исходами} // Математика, механика, 


кибернетика. --


Саратов: Изд-во СГУ, 2000. -- С.\,105--108.


\bibitem{4} 


Бондарева О.Н. {\em Некоторые применения методов линейного 


программирования к теории кооперативных игр} // Пробл. кибернетики. 


-- 1963. -- \No\,10. -- C.\,119--139. 


 


\end{thebibliography} 


 


\vskip 2cm 


 


\post{\it Саратовский государственный университет}{\it Поступила} 


\post{}{02.07.2001} 


 


\label{end} 


\end{document} 


 


\bibitem{5} 


Розен В.В. {\em Вложения упорядоченных множеств в упорядоченные векторные 


пространства} // Упорядоченные множества и решетки. --


Саратов: МВУИП, 1995. -- Вып.\,11. -- C.\,45--52. 


